\documentclass[
    language=english,
    doctype=audit-report,
    institution=none,
]{../../omnilatex}

\RequirePackage{config/document-settings}

\begin{document}

\maketitle

\section*{Independent Auditor's Report}

To the Board of Directors and Shareholders of Meridian Technologies, Inc.:

\subsection*{Opinion}

We have audited the accompanying consolidated financial statements of Meridian
Technologies, Inc. and its subsidiaries, which comprise the consolidated
balance sheet as of December 31, 2025, and the related consolidated
statements of income, comprehensive income, shareholders' equity, and cash
flows for the year then ended, and the related notes to the financial
statements.

In our opinion, the accompanying consolidated financial statements present
fairly, in all material respects, the financial position of Meridian
Technologies, Inc. as of December 31, 2025, and its financial performance
and cash flows for the year then ended in accordance with accounting
principles generally accepted in the United States of America.

\subsection*{Basis for Opinion}

We conducted our audit in accordance with auditing standards generally
accepted in the United States of America. Our responsibilities under those
standards are further described in the Auditor's Responsibilities for the
Audit of the Financial Statements section of our report. We are required to
be independent with respect to the Company in accordance with the ethical
requirements that are relevant to our audit of the financial statements. We
believe that the audit evidence we have obtained is sufficient and appropriate
to provide a basis for our audit opinion.

\section{Scope of the Audit}

The audit was conducted in accordance with generally accepted auditing
standards (GAAS) issued by the American Institute of Certified Public
Accountants. The scope included examination of evidence supporting the
amounts and disclosures in the financial statements, assessment of the
accounting principles used, evaluation of significant estimates made by
management, and consideration of the overall financial statement presentation.

\begin{table}[h]
    \centering
    \begin{tabular}{lll}
        \toprule
        \textbf{Audit Component} & \textbf{Status} & \textbf{Materiality} \\
        \midrule
        Revenue recognition       & Completed & \$2.4M threshold \\
        Inventory valuation       & Completed & \$1.8M threshold \\
        Goodwill impairment       & Completed & \$48M balance \\
        Related-party transactions & Completed & Disclosed \\
        Subsequent events         & Completed & Evaluated \\
        Internal controls (SOX)   & Completed & No material weakness \\
        \bottomrule
    \end{tabular}
    \caption{Summary of audit components and materiality thresholds.}
\end{table}

\section{Consolidated Balance Sheet}

\begin{table}[h]
    \centering
    \begin{tabular}{lrr}
        \toprule
        \textbf{(in thousands)} & \textbf{Dec 31, 2025} & \textbf{Dec 31, 2024} \\
        \midrule
        \textbf{Assets} & & \\
        \quad Cash and equivalents & \$84,200 & \$62,500 \\
        \quad Accounts receivable, net & \$118,400 & \$102,300 \\
        \quad Inventory & \$47,600 & \$41,800 \\
        \quad Prepaid expenses & \$8,900 & \$7,200 \\
        \midrule
        \textbf{Total Current Assets} & \textbf{\$259,100} & \textbf{\$213,800} \\
        \midrule
        \quad Property \& equipment, net & \$96,400 & \$88,100 \\
        \quad Goodwill & \$482,000 & \$482,000 \\
        \quad Intangible assets, net & \$56,300 & \$64,800 \\
        \quad Deferred tax assets & \$12,400 & \$10,900 \\
        \midrule
        \textbf{Total Assets} & \textbf{\$906,200} & \textbf{\$859,600} \\
        \midrule
        \textbf{Liabilities} & & \\
        \quad Accounts payable & \$52,100 & \$46,800 \\
        \quad Accrued liabilities & \$38,600 & \$34,200 \\
        \quad Current portion of long-term debt & \$12,500 & \$12,500 \\
        \midrule
        \textbf{Total Current Liabilities} & \textbf{\$103,200} & \textbf{\$93,500} \\
        \midrule
        \quad Long-term debt & \$225,000 & \$237,500 \\
        \quad Deferred tax liabilities & \$28,400 & \$24,100 \\
        \quad Other non-current liabilities & \$15,800 & \$13,600 \\
        \midrule
        \textbf{Total Liabilities} & \textbf{\$372,400} & \textbf{\$368,700} \\
        \midrule
        \textbf{Shareholders' Equity} & & \\
        \quad Common stock & \$500 & \$500 \\
        \quad Additional paid-in capital & \$312,600 & \$298,400 \\
        \quad Retained earnings & \$225,200 & \$196,000 \\
        \quad Accumulated other comp. loss & \$(4,500) & \$(4,000) \\
        \midrule
        \textbf{Total Shareholders' Equity} & \textbf{\$533,800} & \textbf{\$490,900} \\
        \bottomrule
    \end{tabular}
    \caption{Consolidated balance sheet.}
\end{table}

\section{Consolidated Statement of Income}

\begin{table}[h]
    \centering
    \begin{tabular}{lrrr}
        \toprule
        \textbf{(in thousands)} & \textbf{FY 2025} & \textbf{FY 2024} & \textbf{Change} \\
        \midrule
        \quad Net revenue & \$1,248,600 & \$1,086,200 & +15.0\% \\
        \quad Cost of goods sold & \$(749,160) & \$(662,582) & +13.1\% \\
        \midrule
        \textbf{Gross Profit} & \textbf{\$499,440} & \textbf{\$423,618} & \textbf{+17.9\%} \\
        \textbf{Gross Margin} & \textbf{40.0\%} & \textbf{39.0\%} & \textbf{+1.0pp} \\
        \midrule
        \quad Research \& development & \$(124,860) & \$(108,620) & +15.0\% \\
        \quad Sales \& marketing & \$(187,290) & \$(162,930) & +15.0\% \\
        \quad General \& administrative & \$(74,916) & \$(65,172) & +15.0\% \\
        \midrule
        \textbf{Operating Income} & \textbf{\$212,374} & \textbf{\$186,896} & \textbf{+13.6\%} \\
        \midrule
        \quad Interest expense & \$(11,250) & \$(13,125) & $-$14.3\% \\
        \quad Other income (expense) & \$3,200 & \$1,800 & +77.8\% \\
        \midrule
        \textbf{Income Before Tax} & \textbf{\$204,324} & \textbf{\$175,571} & \textbf{+16.4\%} \\
        \quad Income tax expense & \$(51,081) & \$(43,893) & +16.4\% \\
        \midrule
        \textbf{Net Income} & \textbf{\$153,243} & \textbf{\$131,678} & \textbf{+16.4\%} \\
        \bottomrule
    \end{tabular}
    \caption{Consolidated statement of income.}
\end{table}

\section{Key Audit Matters}

\subsection{Revenue Recognition}

We identified revenue recognition as a key audit matter due to the complexity
of multi-element arrangements and the significant judgments required in
determining standalone selling prices. Our audit procedures included testing
management's methodology for allocating transaction prices, examining a
sample of contracts for proper identification of performance obligations, and
verifying the timing of revenue recognition against delivery and acceptance
records.

\subsection{Goodwill Impairment Testing}

The Company carries \$482 million in goodwill across three reporting units.
Management performs annual impairment testing using a combination of
discounted cash flow analysis and market-based approaches. We evaluated the
reasonableness of key assumptions including discount rates, terminal growth
rates, and revenue projections by comparing them to industry benchmarks and
peer company data.

\begin{table}[h]
    \centering
    \begin{tabular}{lrrr}
        \toprule
        \textbf{Reporting Unit} & \textbf{Goodwill (\$M)} & \textbf{Fair Value (\$M)} & \textbf{Headroom (\%)} \\
        \midrule
        Enterprise Software & \$245 & \$312 & 27.3\% \\
        Cloud Services & \$158 & \$198 & 25.3\% \\
        Professional Services & \$79 & \$104 & 31.6\% \\
        \midrule
        \textbf{Total} & \textbf{\$482} & \textbf{\$614} & \textbf{27.4\%} \\
        \bottomrule
    \end{tabular}
    \caption{Goodwill impairment analysis by reporting unit.}
\end{table}

\section{Internal Controls Over Financial Reporting}

Our audit included obtaining an understanding of internal control over
financial reporting, assessing the risk that a material misstatement exists,
and testing and evaluating the design and operating effectiveness of internal
controls. We did not identify any material weaknesses in internal control
over financial reporting as of December 31, 2025.

\begin{table}[h]
    \centering
    \begin{tabular}{lcc}
        \toprule
        \textbf{Control Area} & \textbf{Effective} & \textbf{Deficiency} \\
        \midrule
        Entity-level controls & Yes & None identified \\
        Revenue \& receivables & Yes & None identified \\
        Inventory \& COGS & Yes & None identified \\
        Payroll \& expenses & Yes & None identified \\
        Financial close process & Yes & None identified \\
        IT general controls & Yes & None identified \\
        \bottomrule
    \end{tabular}
    \caption{Summary of internal control assessment.}
\end{table}

\section{Going Concern}

In connection with our audit of the consolidated financial statements, we
evaluated whether there is substantial doubt about the Company's ability to
continue as a going concern for a reasonable period of time. We concluded
that no such substantial doubt exists. The Company maintains adequate cash
reserves and has access to an undrawn revolving credit facility of
\$150 million.

\section{Management Responsibilities}

The Company's management is responsible for the preparation and fair
presentation of the consolidated financial statements in accordance with
accounting principles generally accepted in the United States of America.
This responsibility includes designing, implementing, and maintaining
internal control relevant to the preparation of financial statements that
are free from material misstatement, whether due to fraud or error.

\begin{description}
    \item[Managing Partner.] Sarah Chen, Deloitte LLP
    \item[Engagement Partner.] Michael Torres, Deloitte LLP
    \item[Date.] March 15, 2026
    \item[Location.] Chicago, Illinois
\end{description}

\end{document}
